\scanpage{38}
\deleted{on the ``pinched'' curve $Y'\deleted{\unsure{=\widehat Y\amalg_I|=\widehat Y/I}}$ (obtained from $\widehat Y$ by contracting $I$ to a point) by identifying the fibers \inserted{reduced over the $y_i\in I$} of the trivial invertible sheaf on $\widehat Y$, \deleted{\unsure{into one}} with $k$ by the $\lambda_i$ -- hence with one another. Then $u+x$ (for $x\in Y$) \inserted{in $J_Y^{\unsure{1|0}}$} is represented by the sheaf $\underline O_{\widehat Y}[x]$ (associated with the divisor $x$), equipped with these $\overline m$s identifications.]}

For \inserted{$x+u$} \deleted{it} to be of the form $y$ $(y\in Y)$ it is already necessary that $u+x$ and $y$ have the $\overline m$ image in $J_{\widehat Y}^1$, \unsure{which|which} \deleted{(for a genus $g\geq1$)} \inserted{and since the image of $u$ in $J_{\widehat Y}^0$ is zero,} already implies \inserted{(for a $g\geq1$)} $x=y$, hence $u=0$. If $g\unsure{=|-}0$, hence $\nu=\operatorname{card}I\geq3$ because of the anabelian hypothesis, we \deleted{\unsure{replace}} take a subset \inserted{$J$\unsure{,}} of $I$ with two elements, and we replace $\widehat Y$ by $\widehat Y\setminus J$ to conclude again that $x=y$ hence $u=0$.

\margin{\unsure{However|How}, we do not know that the image of $u$ in $J_{\widehat Y\setminus J}^0$ is zero!}

[In fact, we have proved that if $u\in\operatorname{Ker}(J_Y^0-J_{\widehat Y}^0)$, then $(Y+u)\cap Y=\unsure{\varnothing|\phi}$ -- and on the other hand, if there exists a nonempty open subset $U$ of $Y$ such that $U+u\subset Y$, then $u$ is in the preceding kernel \unsure{(du m|the latter} $\neq1$). But these arguments work only for \uline{$g\neq1$}.]
\margin{$g\geq2$?}

It is better, in the general case, to present things in \deleted{\unsure{co}}homological form. Let us consider the two \deleted{maps} \inserted{morphisms} $U\overset{i}{\lhook\joinrel\longrightarrow}Y$, inclusion
\margin{37}
